%%% The main file. It contains definitions of basic parameters and includes all other parts.

%% Settings for single-side (simplex) printing
% Margins: left 40mm, right 25mm, top and bottom 25mm
% (but beware, LaTeX adds 1in implicitly)
\documentclass[12pt,a4paper]{report}
\setlength\textwidth{145mm}
\setlength\textheight{247mm}
\setlength\oddsidemargin{15mm}
\setlength\evensidemargin{15mm}
\setlength\topmargin{0mm}
\setlength\headsep{0mm}
\setlength\headheight{0mm}
% \openright makes the following text appear on a right-hand page
\let\openright=\clearpage

%% Settings for two-sided (duplex) printing
% \documentclass[12pt,a4paper,twoside,openright]{report}
% \setlength\textwidth{145mm}
% \setlength\textheight{247mm}
% \setlength\oddsidemargin{14.2mm}
% \setlength\evensidemargin{0mm}
% \setlength\topmargin{0mm}
% \setlength\headsep{0mm}
% \setlength\headheight{0mm}
% \let\openright=\cleardoublepage

%% Generate PDF/A-2u
\usepackage[a-2u]{pdfx}

%% Character encoding: usually latin2, cp1250 or utf8:
\usepackage[utf8]{inputenc}

%% Prefer Latin Modern fonts
\usepackage{lmodern}

%% Further useful packages (included in most LaTeX distributions)
\usepackage{amsmath}        % extensions for typesetting of math
\usepackage{amsfonts}       % math fonts
\usepackage{amsthm}         % theorems, definitions, etc.
\usepackage{amssymb}
\usepackage{bbding}         % various symbols (squares, asterisks, scissors, ...)
\usepackage{bm}             % boldface symbols (\bm)
\usepackage{graphicx}       % embedding of pictures
\usepackage{fancyvrb}       % improved verbatim environment
\usepackage{natbib}         % citation style AUTHOR (YEAR), or AUTHOR [NUMBER]
\usepackage[nottoc]{tocbibind} % makes sure that bibliography and the lists
			    % of figures/tables are included in the table
			    % of contents
\usepackage{dcolumn}        % improved alignment of table columns
\usepackage{booktabs}       % improved horizontal lines in tables
\usepackage{paralist}       % improved enumerate and itemize
\usepackage{xcolor}         % typesetting in color

%%% My packages

\usepackage{stmaryrd}
\usepackage{mathtools}

%%% Basic information on the thesis

% Thesis title in English (exactly as in the formal assignment)
\def\ThesisTitle{Pseudofinite structures}

% Author of the thesis
\def\ThesisAuthor{Ondřej Ježil}

% Year when the thesis is submitted
\def\YearSubmitted{2022}

% Name of the department or institute, where the work was officially assigned
% (according to the Organizational Structure of MFF UK in English,
% or a full name of a department outside MFF)
\def\Department{Department of Algebra}

% Is it a department (katedra), or an institute (ústav)?
\def\DeptType{Department}

% Thesis supervisor: name, surname and titles
\def\Supervisor{prof. RNDr. Jan Krajíček Ph.D.}

% Supervisor's department (again according to Organizational structure of MFF)
\def\SupervisorsDepartment{Department of Algebra}

% Study programme and specialization
\def\StudyProgramme{Matemathics}
\def\StudyBranch{Mathemathics for Information technologies}

% An optional dedication: you can thank whomever you wish (your supervisor,
% consultant, a person who lent the software, etc.)
\def\Dedication{%
Dedication.
}

% Abstract (recommended length around 80-200 words; this is not a copy of your thesis assignment!)
\def\Abstract{%
Abstract.
}

% 3 to 5 keywords (recommended), each enclosed in curly braces
\def\Keywords{%
{key} {words}
}

%% The hyperref package for clickable links in PDF and also for storing
%% metadata to PDF (including the table of contents).
%% Most settings are pre-set by the pdfx package.
\hypersetup{unicode}
\hypersetup{breaklinks=true}

% Definitions of macros (see description inside)
%%% This file contains definitions of various useful macros and environments %%%
%%% Please add more macros here instead of cluttering other files with them. %%%

%%% Minor tweaks of style

% These macros employ a little dirty trick to convince LaTeX to typeset
% chapter headings sanely, without lots of empty space above them.
% Feel free to ignore.
\makeatletter
\def\@makechapterhead#1{
  {\parindent \z@ \raggedright \normalfont
   \Huge\bfseries \thechapter. #1
   \par\nobreak
   \vskip 20\p@
}}
\def\@makeschapterhead#1{
  {\parindent \z@ \raggedright \normalfont
   \Huge\bfseries #1
   \par\nobreak
   \vskip 20\p@
}}
\makeatother

% This macro defines a chapter, which is not numbered, but is included
% in the table of contents.
\def\chapwithtoc#1{
\chapter*{#1}
\addcontentsline{toc}{chapter}{#1}
}

% Draw black "slugs" whenever a line overflows, so that we can spot it easily.
\overfullrule=1mm

%%% Macros for definitions, theorems, claims, examples, ... (requires amsthm package)

\theoremstyle{plain}
\newtheorem{thm}{Theorem}
\newtheorem{lemma}[thm]{Lemma}
\newtheorem{claim}[thm]{Claim}

\theoremstyle{plain}
\newtheorem{defn}{Definition}

\theoremstyle{remark}
\newtheorem*{cor}{Corollary}
\newtheorem*{rem}{Remark}
\newtheorem*{example}{Example}

%%% An environment for proofs

\newenvironment{myproof}{
  \par\medskip\noindent
  \textit{Proof}.
}{
\newline
\rightline{$\qedsymbol$}
}

%%% An environment for typesetting of program code and input/output
%%% of programs. (Requires the fancyvrb package -- fancy verbatim.)

\DefineVerbatimEnvironment{code}{Verbatim}{fontsize=\small, frame=single}

%%% The field of all real and natural numbers
\newcommand{\R}{\mathbb{R}}
\newcommand{\N}{\mathbb{N}}

%%% Useful operators for statistics and probability
\DeclareMathOperator{\pr}{\textsf{P}}
\DeclareMathOperator{\E}{\textsf{E}\,}
\DeclareMathOperator{\var}{\textrm{var}}
\DeclareMathOperator{\sd}{\textrm{sd}}

%%% Transposition of a vector/matrix
\newcommand{\T}[1]{#1^\top}

%%% Various math goodies
\newcommand{\goto}{\rightarrow}
\newcommand{\gotop}{\stackrel{P}{\longrightarrow}}
\newcommand{\maon}[1]{o(n^{#1})}
\newcommand{\abs}[1]{\left|{#1}\right|}
\newcommand{\dint}{\int_0^\tau\!\!\int_0^\tau}
\newcommand{\isqr}[1]{\frac{1}{\sqrt{#1}}}

%%% Various table goodies
\newcommand{\pulrad}[1]{\raisebox{1.5ex}[0pt]{#1}}
\newcommand{\mc}[1]{\multicolumn{1}{c}{#1}}

%%% Custom macros
\newcommand{\Th}{\text{Th}}
\newcommand{\st}{\text{st}}
\newcommand{\bbl}{\llbracket}
\newcommand{\bbr}{\rrbracket}
\newcommand{\Lor}{\bigvee}
\newcommand{\Land}{\bigwedge}
\newcommand{\1}{\textbf{1}}
\newcommand{\0}{\textbf{0}}
\newcommand{\dpt}{\text{dp}}
\newcommand{\Aut}{\text{Aut}}
\newcommand{\pPATH}{*\text{PATH}}
\newcommand{\EDGE}{\text{EDGE}}
\newcommand{\nEDGE}{\text{nonEDGE}}
\newcommand{\ALL}{\text{ALL}}
\newcommand{\leqp}{\trianglelefteq}
\newcommand{\set}[1]{\langle #1)}
\newcommand{\ceil}[1]{\left \lceil {#1} \right \rceil}
\newcommand{\floor}[1]{\left \lfloor {#1} \right \rfloor}
\newcommand{\M}{\mathcal{M}}
\newcommand{\Po}{\mathcal{P}}
\newcommand{\A}{\mathcal{A}}
\newcommand{\B}{\mathcal{B}}
\newcommand{\I}{\mathcal{I}}
\newcommand{\card}{\text{card}}
\newcommand{\SK}{\text{SK}}
\newcommand{\CK}{\text{CK}}
\newcommand{\G}{\mathcal{G}}
\newcommand{\Q}{\mathbb{Q}}
\renewcommand{\P}{\textbf{P}}
\newcommand{\NP}{\textbf{NP}}
\newcommand{\FP}{\textbf{FP}}
\newcommand{\TFNP}{\textbf{TFNP}}
\newcommand{\LEAF}{\text{LEAF}}
\newcommand{\SOURCEORSINK}{\text{SOURCE.OR.SINK}}
\newcommand{\PPA}{\textbf{PPA}}
\newcommand{\PPAD}{\textbf{PPAD}}
\renewcommand{\N}{\mathbb{N}}
\newcommand{\dom}{\text{dom}\:}
\newcommand{\Rng}{\text{Rng}\:}
\newcommand{\Frac}{\text{Frac}}
%\newcommand{\abs}[1]{\lvert #1 \rvert}


% Definitions of amsthm environments
\newtheoremstyle{mystyle}%                % Name
  {}%                                     % Space above
  {}%                                     % Space below
  {}%                             % Body font
  {}%                                     % Indent amount
  {\bfseries}%                            % Theorem head font
  {.}%                                    % Punctuation after theorem head
  { }%                                    % Space after theorem head, ' ', or \newline
  {}%                                     % Theorem head spec (can be left empty, meaning `normal')

\theoremstyle{mystyle}

\newtheorem{thrm}{Theorem}[section]
\newtheorem{lemm}[thrm]{Lemma}
\newtheorem{exam}[thrm]{Example}
\newtheorem{defi}[thrm]{Definition}
\newtheorem{crll}[thrm]{Corollary}

% Title page and various mandatory informational pages
\begin{document}
%%% Title page of the thesis and other mandatory pages

%%% Title page of the thesis

\pagestyle{empty}
\hypersetup{pageanchor=false}
\begin{center}

\centerline{\mbox{\includegraphics[width=166mm]{../img/logo-en.pdf}}}

\vspace{-8mm}
\vfill

{\bf\Large MASTER THESIS}

\vfill

{\LARGE\ThesisAuthor}

\vspace{15mm}

{\LARGE\bfseries\ThesisTitle}\footnote{This is a second version of a thesis originally defended in June 2022. It includes errata and several parts were rewritten to improve clarity.}

\vfill

\Department

\vfill

{
\centerline{\vbox{\halign{\hbox to 0.45\hsize{\hfil #}&\hskip 0.5em\parbox[t]{0.45\hsize}{\raggedright #}\cr
Supervisor of the master thesis:&\Supervisor \cr
\noalign{\vspace{2mm}}
Study programme:&\StudyProgramme \cr
\noalign{\vspace{2mm}}
Study branch:&\StudyBranch \cr
}}}}

\vfill

% Zde doplňte rok
Prague \YearSubmitted

\end{center}

\newpage

%%% Here should be a bound sheet included -- a signed copy of the "master
%%% thesis assignment". This assignment is NOT a part of the electronic
%%% version of the thesis. DO NOT SCAN.

%%% A page with a solemn declaration to the master thesis

%%%\openright
%%%\hypersetup{pageanchor=true}
%%%\pagestyle{plain}
%%%\pagenumbering{roman}
%%%\vglue 0pt plus 1fill
%%%
%%%\noindent
%%%I declare that I carried out this master thesis independently, and only with the cited
%%%sources, literature and other professional sources. It has not been used to obtain another
%%%or the same degree.
%%%
%%%\medskip\noindent
%%%I understand that my work relates to the rights and obligations under the Act No.~121/2000 Sb.,
%%%the Copyright Act, as amended, in particular the fact that the Charles
%%%University has the right to conclude a license agreement on the use of this
%%%work as a school work pursuant to Section 60 subsection 1 of the Copyright~Act.
%%%
%%%\vspace{10mm}
%%%
%%%\hbox{\hbox to 0.5\hsize{%
%%%In \hbox to 6em{\dotfill} date \hbox to 6em{\dotfill}
%%%\hss}\hbox to 0.5\hsize{\dotfill\quad}}
%%%\smallskip
%%%\hbox{\hbox to 0.5\hsize{}\hbox to 0.5\hsize{\hfil Author's signature\hfil}}
%%%
%%%\vspace{20mm}
%%%\newpage
%%%
%%%%%% Dedication
%%%
%%%\openright
%%%
%%%\noindent
%%%\Dedication
%%%
%%%\newpage

%%% Mandatory information page of the thesis

\openright

\vbox to 0.5\vsize{
\setlength\parindent{0mm}
\setlength\parskip{5mm}

Title:
\ThesisTitle

Author:
\ThesisAuthor

\DeptType:
\Department

Supervisor:
\Supervisor, \SupervisorsDepartment

Abstract:
\Abstract

Keywords:
\Keywords

\vss}

\newpage

\openright
\pagestyle{plain}
\pagenumbering{arabic}
\setcounter{page}{1}


%%% A page with automatically generated table of contents of the master thesis

\tableofcontents

%%% Each chapter is kept in a separate file
\chapter*{Introduction}
\addcontentsline{toc}{chapter}{Introduction}

There exist several logical constructions of limits of classes of finite structures such as the ultraproduct construction or using the compactness theorem. The latter used in \cite{Fagin1976} to prove the 0-1 law for structures over relational vocabularies.

In combinatorics there are also several notions of limits of finite graphs. For example the dense graph limit defined for a sequence of graph $\{G_k\}_{k>0}$ satisfying the condition that
\[t(F,G_n)=\frac{\hom(F,G)}{\abs{G_n}^{\abs{F}}}\]
converges for every fixed connected graph $F$ which provided a framework (see \cite{lovasz2006limits}) to restate and find new proofs for results in extremal graph theory such as Goodman's theorem relating the number of edges to the number of triangles in a graph. There are other notions of limits of sequences of graphs we recommend to read \cite{Nesetril2013} to the interested reader. Another recent use of limit objects for the results of extremal combinatorics was by Razborov in \cite{razborov2007flag}.

These diferent notions of limits directly or tangentially relate to the concept of pseudofinite structures. For a first order language $L$, we call an $L$-structure $S$ pseudofinite, if it satisfies the theory $T_f$ consisting of all sentences true in all finite $L$-structures.

In this thesis we use the concept of pseudofinite structures to define a limit of a family of finite graphs relative to some computationally restricted class of functions. Instead of studying the density of substructures we study these so called wide limits both generally and by analyzing concrete examples and tying them with complexity-theoretical statements.

The key method we use is arithmetical forcing with random variables which allows us to construct models of (weak) arithmetical theories and by restricting to a language of graphs gives us boolean valued graphs. An important resource for this technique is \cite{krajicek2010forcing}. In these boolean valued graphs, witnessing of existential quantifiers corresponds to complexity theoretical properties of the class of graphs we are considering. 

(TODO: Comment on structure of chapters.)

\chapter{The limit}
\section{Basic examples}
\subsection{$\mathcal{G}_k=\EDGE_k$}
We consider the classes of graphs
\[\EDGE_k:=\{\text{$(k,E); \abs{E}=1$}\},\]
and we let $\mathcal{G}_k=\EDGE_k$ and $F=F_{tree}$.

Intuitivelly one should not be able to find the edge on a significant i.e. non-infinitesimal fraction of samples with a tree that is allowed to explore only an infinitesimal fraction of edges.

\begin{thrm}\label{thrmedge}
We will prove that \[\bbl(\exists x)(\exists y)\Gamma(x,y)\bbr=\textbf{0}.\]
\end{thrm}
\begin{proof}
Let $T\in F_{tree}$ be a tree of depth $n^{1/t}$, for some $t>\mathbb{N}$ that outputs a pair of numbers less than $n$.

Start from the root of $T$ and always choose the path that corresponds to an edge not existing. At the end we obtain some answer, that gives us a set of at most $2 \cdot n^{1/t}+2$ vertices. Now we can find at least: 
\begin{align}
\binom{n-2n^{1/t}-2}{2}&=\frac{(n-2n^{1/t}-2)(n-2n^{1/t}-3)}{2}
\end{align}
different $\omega\in\Omega$ such that $T(\omega)$ is not an edge in $\omega$. 

The probability that any of those graphs is sampled is
\begin{align}
\frac{\binom{n-2n^{1/t}-2}{2}}{\abs{\mathcal{G}_k}}&=\frac{(n-2n^{1/t}-2)(n-2n^{1/t}-3)}{n(n-1)}\\
&=\left(1-\frac{2n^{1/t}-2}{n}\right)\cdot\left(1-\frac{2n^{1/t}-3}{n-1}\right)\\
&\geq \left(1-\frac{2n^{1/t}-2}{n}\right)^2\\
&\geq 1-\frac{4n^{1/t}-4}{n}.
\end{align}
And one can clearly see that $\st(1-\frac{4n^{1/t}-4}{n})=1$. This proves that the boolean value we are considering is $\textbf{0}$ since we can the two witnesses for $x$ and $y$ into a tree that could find an edge with depth $n^{1/t}$ for some $t>\mathbb{N}$.
\end{proof}

\subsection{Sparse $\mathcal{G}_k$}

One can see that in Theorem \ref{thrmedge} the exact shape of graphs in $\mathcal{G}_k$ does not play a crucial role. If $\mathcal{G}_k$ consisted of all graphs on $k$ verticies containing say exactly one triangle, or any other fixed subgraph of constant size, and no other edges, we would still find that the non-existence is valid in the limit graph.

A more general case would be to consider a family of graphs in which there is an infinitesimally small chance that two independent uniformly random verticies have an edge between. However, this is not sufficient.

\begin{exam}
Let $\mathcal{G}_k$ consist of those graphs on the vertex-set $k$ which contain the edge $E(0,1)$ and then has exactly one other edge. 

As $k$ increases, the number of edges get smaller than any standard positive fraction. But
\[\bbl(\exists x)(\exists y)E(x,y)\bbr=\1,\]
as witnessed by $x$ being the constant $0$ and $y$ the constant $1$ both of which are computable by a tree of depth $0$.
\end{exam}

One can see that having distinguished verticies can ruin the sparseness implying the non-existence of edges in the limit graph. We want to distiguish from this situation by considering the sequences $\mathcal{G}_k$ satisfying the following natural definition.

\begin{defi}
We say that $\{\mathcal{G}_k\}_{k=0}^\infty$ is \textbf{isomorphism closed}, if there is $k_0$ such that for every $k>k_0$ if we have that $G_1\in\mathcal{G}_k$, $V_{G_2}=g_k$ and $G_1\cong G_2$ then $G_2\in\mathcal{G}_k$.
\end{defi}


\begin{thrm}\label{thrmsparse}
Let an isomorphism closed $\mathcal{G}_k$ have the following property. There is a sequence $\{b_k\}_{k}$ and for big enough $k$, a uniformly sampled $2$-element $\{u,v\}\subseteq g_k$ and every $G\in\mathcal{G}_k$ we have
\[\Pr[E_G(u,v)]\leq b_k,\]
and some $k_0$ such that $\lim_{k\to\infty}k^{1/k_0}b_k = 0.$ Then

\[\bbl(\exists x)(\exists y)\Gamma(x,y)\bbr=\0.\]
\end{thrm}
\begin{proof}
Let us define the number $c_{u,v}:=\abs{\{G\in\mathcal{G}_k; E_G(u,v)\}}$, which is the number of graphs $G$ in $\mathcal{G}_k$ satisfying $E_G(u,v)$. Of course $c_{u,u}=0$ for every $u$.

\vspace{0.5em}
\textbf{Claim:} Let $u\neq v,u'\neq v'$ be vertices, then $c_{u,v}=c_{u',v'}$.\\
\textit{proof of claim:} Let $\rho:=(u\:u')(v\:v')$ be a permutation with cycles $(u\:u')$ and $(v\:v')$. We can let $\rho$ act on $\mathcal{G}_k$ by sending $G$ to a graph $\rho(G)$ which renames the edges coordinate-wise.

The fact that $\mathcal{G}_k$ is isomorphism closed implies that $\rho$ restricts to a bijection:
\[\rho':\{G\in\mathcal{G}_k;E_G(u,v), \lnot E_G(u',v')\}\to \{G\in\mathcal{G}_k;E_G(u',v'),\lnot E_G(u,v)\}\]
which proves the claim. \qed
\vspace{0.5em}

Now we define a matrix with entries
\[a_{G,\{u,v\}}:=\begin{cases}1&E_G(u,v)\\0&\text{otherwise}\end{cases}\]
where the rows are indexed by one of $\abs{\mathcal{G}_k}$ many graphs in $\mathcal{G}_k$ and the columns are indexed by the $\frac{k(k-1)}{2}$ many $2$-element sets of numbers in $k$. We take any distinct vertices $u,v$ and define $p:=\Pr_{G\in \mathcal{G}_k}[E_G(u,v)]=\frac{c_{u,v}}{\abs{\mathcal{G}_k}},$ by the claim the choice of $u,v$ does not matter.

The assumption from the statement is equivalent to the equality 
\[\sum_{\{u,v\}}a_{G,\{u,v\}}\leq \frac{k(k-1)}{2}b_k\] 
for every $G$. We combine this with the claim and the definition of $p$ to get
\begin{align}
\frac{k(k-1)}{2}\abs{\mathcal{G}_k}p&=\sum_{\{u,v\}}\sum_{G\in\mathcal{G}_k}a_{G,\{u,v\}}\\
&=\sum_{G\in\mathcal{G}_k}\sum_{\{u,v\}}a_{G,\{u,v\}}\\
&\leq \abs{\mathcal{G}_k}\frac{k(k-1)}{2}b_k
\end{align}
which implies
\[p\leq b_k.\]

Now let $k:=n$ and let $T\in F_{tree}$ be a tree of depth $n^{1/t}$ for some $t>\mathbb{N}$, where every leaf of $T$ is labeled by some edge. Walk down the tree $T$ from the root by answering negatively to every edge, this gives us a set $E_T$ of all edges $T$ inspected or outputed and $\abs{E_T}\leq n^{1/t}+1$.

Now we just need to prove that the probability $T$ find an edge is infinitesimally small. This is enough to prove the theorem, since the trees computing any two witnesses for $x$ and $y$ in the statement can be combined to construct $T$ and if any tree $T$ succeeds with only infinitesimally small probability, no random vertices can witness an edge on a set of non-zero measure.

We use the fact that $p\leq b_n$ to derive 
\begin{align}
\Pr_{G\in\mathcal{G}_n}[T\text{ finds an edge}]&\leq \sum_{\{u,v\}\in E_T} \Pr_{G\in \mathcal{G}_n}[E_G(u,v)]\\
&=\sum_{\{u,v\}\in E_T} \frac{c_{u,v}}{\abs{\mathcal{G}_n}}\\
&\leq\sum_{\{u,v\}\in E_T} p\\
&=(n^{1/t}+1)p\\
&\leq(n^{1/t}+1)b_k\\
&\leq n^{1/k_0}b_k\\
&\approx 0,
\end{align}
which proves the theorem.
\end{proof}

The assumption $\lim_{k\to \infty} k^{1/k_0}b_k=0$ for some $k_0$ may seem unintuitive at first. However, it precisely reflects what is ``sparse" for the trees in $T_{tree}$. The following example shows that without the assumption the theorem fails.

\begin{exam}
Let $\mathcal{G}_k$ consist of all graphs on the vertex set $\set{k}$ with precisely $\lceil\frac{k(k-1)}{2\log k}\rceil$ edges.

Then we claim that
\[\bbl(\exists x)(\exists y)\Gamma(x,y)\bbr=\1.\]

Let $\alpha$ and $\beta$ be verticies computed by the tree of the same shape which inspects a set of any $n^{1/t}$ distinct edges for some $t>\mathbb{N}$. If it finds an edge we define $\alpha$ and $\beta$ in any way so they are the distinct verticies incidental with this edge. Otherwise we let $\alpha(\omega)=\beta(\omega)=0$.

Let $T$ be a tree of the same shape, that computes the pair $\{\alpha,\beta\}$ then we can compute the probability where such a tree fails as the fraction of all graphs which have no edges that $T$ inspects. Let $m=\binom{n}{2}$. We get

\begin{align}
\Pr_{G\in\mathbb{G}_n}[T\text{ fails}]&=\frac{\binom{m-n^{1/t}}{\ceil{ \frac{n(n-1)}{2\log n}} }}{\binom{m}{\ceil{ \frac{n(n-1)}{2\log n}}}}\\
&=\frac{\frac{(m-n^{1/t})!}{\ceil{\frac{n(n-1)}{2\log n}} ! (m-\ceil{ \frac{n(n-1)}{2\log n}}-n^{1/t})!}}{\frac{m!}{\lceil \frac{n(n-1)}{2\log n}\rceil!(m-\ceil{\frac{n(n-1)}{2\log n}})!}}\\
&=\frac{(m-n^{1/t})!(m-\ceil{ \frac{n(n-1)}{2\log n}})!}{m!(m-\ceil{\frac{n(n-1)}{2\log n}}-n^{1/t})!}\\
&=\prod_{i=0}^{n^{1/t}-1}\frac{m-\ceil{\frac{n(n-1)}{2}}-i}{m-i}\\
&=\left(1-\frac{\ceil{\frac{n(n-1)}{2\log n}}}{\frac{n(n-1)}{2}}\right)^{n^{1/t}}\\
&\leq\left(1-\frac{\ceil{\frac{n(n-1)}{2\log n}}}{\frac{n(n-1)}{2}}\right)^{n^{1/t}}\\
&\leq\left(1-\frac{1}{\log n}\right)^{n^{1/t}}
\end{align}

And for any standard $k$ we have
\begin{align}
\left(1-\frac{1}{\log n}\right)^{n^{1/t}}&\leq\left(1-\frac{1}{\log n}\right)^{k\cdot \log n}\\
&\leq (e^{-\frac{1}{\ln 2}})^k.
\end{align}

So $\st(\Pr_{G\in\mathcal{G}_n}[T\text{ fails}])=0$ and we get 
\begin{align}
\mu(\bbl(\exists x)(\exists y)\Gamma(x,y)) \bbr &\geq \mu(\bbl \Gamma(\alpha,\beta)\bbr)\\
&=\st(1-\Pr_{G\in\mathbb{G}_n}[T\text{ fails}])\\
&=1.
\end{align}
\end{exam}

\section{Isomorphism closed $\mathcal{G}_k$}

So far the measure of every truth value we encountered was either $0$ or $1$. Is there a sequence $\mathcal{G}_k$ whose $F_{tree}$-limit and an $\{E\}$-sentence $\varphi$ such that ${0<\bbl\varphi^\Gamma\bbr<1?}$ As in the case of edge existance for a limit of sparse graphs, it is not hard to come up with an example if we allow distiguishing elements in $\mathcal{G}_k$.

\begin{exam}
Let \[\mathcal{G}_k=\{(\set{k},E);\text{$E$ has exactly two edges, one of them being $\{0,1\}$}\},\]

and let $F=F_{tree}$. Then $\mu(\bbl(\exists x)(\exists y)\Gamma(x,y)\bbr)=\frac{1}{2}$.
\end{exam}
\begin{proof}
Let $T_0$ be a tree that always outputs $0$ and $T_1$ be a tree that always outputs $1$. We can prove that $\bbl\Gamma(0,1)\bbr\geq\bbl\Gamma(\alpha,\beta)\bbr$ for any $\alpha,\beta$. 
\end{proof}

For the case of isomorphism closed $\mathcal{G}_k$ we prove that every $\{E\}$-sentence has truth value either $\0$ or $\1$. We start with the existential case.

\begin{exam}
NOT TRUE! Consider all graphs with one edge, or one non-edge. $\mu(\bbl(\exists x)(\exists y)\Gamma(x,y)\bbr)=\frac{1}{2}.$
\end{exam}

\begin{thrm}
Let $\varphi_0(\overline x)$ be an open $\{E\}$-formula, and let $\mathcal{G}_k$ be isomorphism closed and categorical then
\[\bbl(\exists \overline x)\varphi_0^\Gamma(\overline x)\bbr\in\{\0,\1\}.\]
\end{thrm}

\begin{proof}
Let $\overline T$ be a tuple of trees computing $\overline \alpha$ such that 
\[p:=\mu(\bbl\varphi_0(\overline \alpha)\bbr)>\0.\]

We want to iterate $\overline T$ to amplify the probability of success. Let 
\[A=\{\omega \in \Omega; \omega \models \varphi_0(\overline\alpha(\omega)) \}\]
and $\pi\in S_{\abs{\G_n}}$ with maximal $\abs{\pi[A]\setminus A}$. Let $\pi(\overline T)$ be a tuple of relabeled trees computing $\pi(\overline \alpha)$.
\end{proof}

\chapter{Dense case}\label{chapdense}

\section{$\G_k=\SK_k^{1/2}$}

Now we turn to analyze dense wide sequences in the second order case; that is the case when we shall be interested in second order properties of the wide limit. In this chapter we assume $F=F_{rud}$ and $G=G_{rud}$. Specifically we will consider the problem of finding a large clique in a graph. 

Generally it is considered a computationally hard problem to find a large clique in a graph. It is an $\NP$-hard problem and thus it is conjectured that it cannot be solved in polynomial time. We first turn to the following wide sequence.

\begin{defi}
Let \[\SK_k^{1/2}=\{(\set{k},E);\text{$E$ consists of exactly one $k/2$-clique}\}.\]
\end{defi}

Limiting inputs to $\SK_k^{1/2}$ makes the problem less complex, because for a vertex $v$ to be a part of the biggest clique it is enough that it has nonzero degree. Naturally we want to see whether there is a large clique in $\lim_F^G \SK_n^{1/2}$, every sample has a clique of size $n/2$, but is there a set of $F$-vertices witnessing that? 

Here to measure the size of such a clique, the second order wide limit  by itself is not a sufficient object. Instead we need to turn to the arithmetical expansion $K(\G_n,F,G)$ and ask whether we can find an injective function from some large initial segment into a clique. This can be expressed as a second order sentence with parameters from $F$. It is not hard to prove the following result which also demonstrates how the sentence is formed.

\begin{thrm}\label{thrmsk1}For every $t>\N$ we have
\[K(\SK_k^{1/2},F,G)\bbl\text{$E$ has a clique of size $n^{1/t}$}\bbr=\1.\]

\end{thrm}

\begin{proof}(sketch) We need to analyze the value
\begin{align}
&\bbl(\exists \Lambda)(\forall u)(\forall v)(((u\not = v \leq n^{1/t})\to E(\Lambda(u),\Lambda(v)))\land(\Lambda:[n^{1/t}]\hookrightarrow\mathcal{M}))\bbr
\end{align}
which is equal to
\begin{align}
\Lor_\Lambda\Land_u\Land_v\bbl(u\not=v < n^{1/t})\to(E(\Lambda(u),\Lambda(v))\land\Lambda(u)\not=\Lambda(v))\bbr.
\end{align}
So we want to find some $n^{1/t}$-tuple of trees computing some $\Lambda$, computed by some $(\gamma_0,\dots,\gamma_{n^{1/t}-1})$, which is injective on $\set{n^{1/t}}$ and its $\set{n^{1/t}}$ range is a clique in $E$.

We define $(\gamma_0,\dots,\gamma_{n^{1/t}-1})$ as follows. The tree $T_0$ computing $\Lambda_0$ inspects all the edges $(u,v)\in\set{n^{1/t}}\times\set{n^{1/t}}$ in some specified order and outputs the first vertex it finds with an edge and outputs $0$ otherwise. The tree $T_i$ computing $\gamma_i$ extends $T_{i-1}$ by not outputting the $i$-th vertex with the property but instead continues inspecting and outputs the $(i+1)$-th vertex for which it finds an edge. Also every such tree has depth at most $n^{1/t}\cdot n^{1/t}=n^{2/t}=n^{1/(t/2)}$.

One can check that the probability the tree $\gamma_i$ does not find $i$ vertices with an edge is infinitesimal and therefore it always outputs a vertex in the clique of $\omega$. Intuitively this is because the expected ratio of the vertices in $\set{n^{1/t}}$ which are also in the clique is about one half. Moreover, every $\gamma_i$ outputs the $i$-th element element of the ordering and thus it is injective.
\end{proof}

At first glance the lower bound $n^{1/t}$ for every nonstandard $t$ may seem optimal given the proof method we used, but there is a way to radically improve it. The idea is to partition the set of vertices into many smaller ones and let $\gamma_i$ search only in the $i$-th set. First we need the following lemmas.

\begin{lemm}\label{lemmmiss}
Let $S\subseteq \set{n}$ such that $\abs{S}=m > \N$, then
\[\st\left(\Pr_{G\in\G_n}[\text{$S$ contains no vertices in the clique of $G$}]\right)=0.\]
\end{lemm}
\begin{proof}
There are $\binom{n-m}{\floor{\frac{n}{2}}}$ different graphs in $\SK_n^{1/2}$ in which the clique does not intersect $S$. We then bound the probability as
\begin{align}
\frac{\binom{n-m}{\floor{\frac{n}{2}}}}{\binom{n}{\floor{\frac{n}{2}}}}&=\frac{(n-m)!(n-\floor{\frac{n}{2}})!}{(n)!(n-\floor{\frac{n}{2}}-c)!}\\
&=\prod_{i=0}^{m-1}\frac{(n-i-\floor{\frac{n}{2}})}{(n-i)}\\
&=\prod_{i=0}^{m-1}\left(1-\frac{\floor{\frac{n}{2}}}{n-i}\right)\\
&\leq \left(1-\frac{\floor{\frac{n}{2}}}{n}\right)^m\\
&\leq \left(1-\frac{\floor{\frac{n}{2}}}{n}\right)^{n \cdot \frac{m}{n}}\\
&\leq e^{-\floor{\frac{n}{2}}\frac{m}{n}}.
\end{align}
But $\floor{\frac{n}{2}} \frac{m}{n}$ is infinite therefore the bound is infinitesimal.
\end{proof}

\begin{lemm}\label{lemmfindtree}
Let $a\in\M$, let $v_0,\dots,v_{a-1}\in \set{n}$ distinct vertices, then there exist trees $T_{v_i}$ such that
\[\st\left(\Pr_{\omega\in\Omega}[\forall i: (v_i,T_{v_i}(\omega))\in E_\omega|\forall i: \text{$v_i$ is in the clique}]\right)=1.\]
\end{lemm}

\begin{proof}
The tree $T_{v_i}$ inspects all the edges $(v_i,j)$ where $j$ ranges over $\set{n^{1/t}}$ for some $t>\N$ and outputs $j$ if $(v_i,j)\in E_\omega$. By Lemma \ref{lemmmiss} we have that only an infinitesimal number of graphs have their clique not intersect $\set{n^{1/t}}$ so each $T_{v_i}$ succeeds on all but an infinitesimally small portion of $\Omega$. But if one $T_{v_i}$ finds a neighbour of $v_i$, then all do, since nonzero degree vertices in every $\omega$ form a clique and the same $w\in\set{n^{1/t}}$ is a neighbour of all $v_i$s.
\end{proof}

\begin{lemm}\label{lemmmiss2}
Let $S_0,\dots,S_{a-1}\subseteq\set{n}$ sets of size $a\in \M$ for $i\in \M$ then
\[\Pr_{\omega\in\Omega}\left[\bigcup_{i=0}^{a-1}[\text{$S_i$ contains no vertices in the clique of $\omega$}]\right]\leq a\cdot e^{-\floor{\frac{n}{2}}\frac{m}{n}}.\]
\end{lemm}
\begin{proof}Follows from the proof of Lemma \ref{lemmmiss} and union bound.
\end{proof}

Now we are ready to improve on Theorem \ref{thrmsk1}.

\begin{thrm}\label{thrmsk}
Let $m = 2 \ln n$. Then
\[K(\SK_k^{1/2},F,G)\bbl\text{$E$ has a clique of size $\floor{n/m}$}\bbr=\1.\]
\end{thrm}
\begin{proof}
Partition a subset of $\set{n}$ to sets $S_{0},\dots,S_{\floor{n/m}-1}$ each of size at least $m$. Specifically if $m$ divides $n$ then we partition the whole $\set{n}$.

By Lemma \ref{lemmmiss2} we have that with probability that we do not sample $\omega$ which have the clique intersect all $S_i$s 
\begin{align}
\floor{\frac{n}{m}}\cdot e^{-\floor{\frac{n}{2}}\frac{m}{n}} = e^{\ln\floor{\frac{n}{m}}-\floor{\frac{n}{2}}\frac{m}{n},}\label{eqpr}
\end{align}
we can bound the exponent as
\begin{align}
\ln\floor{\frac{n}{m}}-\floor{\frac{n}{2}}\frac{m}{n}&\leq\ln\frac{n}{m}-\frac{n}{2}\cdot\frac{m}{n}+\frac{m}{n}\\
&\leq\ln n - \ln m - \frac{m}{2} + \frac{m}{n}1\\
&\leq \ln n - \ln \ln n - \ln n+\frac{2\ln n}{n}\\
&\leq - \ln \ln n + \frac{2\ln n}{n}
\end{align}
which is negative and infinite, because $\frac{\ln x}{x}\overset{x\to \infty}{\to} 0$ as  and therefore \eqref{eqpr} is infinitesimal.

So with probability infinitesimally close to 1 we have in each $S_i$ a vertex $v_i$ which is also a part of the clique. By Lemma \ref{lemmfindtree} we have that there exists a tree verifying whether a given vertex is in the clique and since $m\leq n^{1/t}$ for some $t$ we can concatenate the trees to get a tree $T_{S_i}$ which finds in $S_i$ an element of the clique with probability infinitesimally close to $1$.

Finally we can have a function $\Lambda\in G$ computed by $(\Lambda_0,\dots,\Lambda_{\floor{n/m}-1})$ by letting $\Lambda_i$ be computed by $T_{S_i}$ we have already verified $\bbl\Lambda\text{ is a clique}\bbr=\1.$

Because $T_{S_i}(\omega)\in S_i$ when $S_i$ succeeds, and $S_i$ are disjoint we have \[\bbl\Lambda:\set{\floor{n/m}}\xhookrightarrow{} \M\bbr=\1. \]
Which proves the theorem.
\end{proof}

One can also check that for $\SK_{k}^{1/l}$, the graphs whose edges are exactly one $\floor{k/l}$ clique for some $l\in \N$, the wide limit has a clique of size $\floor{n/m}$ for $m=l\cdot \ln(m)$ by the same technique.

Even though the size of the clique has radically increased, we still did not find a clique in $E$ of size $n/2$. One can verify that with the method provided one cannot obtain such a clique because the probability that any of $\floor{n/2}$ two-element sets does not intersect the clique is too large. 

\begin{ques}
Is $\floor{n/(2\ln n)}$ the clique number in the Boolean valued graph $\lim_{F_{rud}}^{G_{rud}}\SK_n^{1/2}$, measured in $K(\SK_k^{1/2},F,G)$?
\end{ques}


\section{$\G_k=\CK_k^{1/2}$}

Now let us mention the more complex case of the wide sequence $\CK_k^{1/2}$
\[\CK_k^{1/2}=\{(\set{k},E);E\text{ contains a $\floor{k/2}$ clique.}\}\]

We are still guaranteed that every $\omega$ contains a large clique, but there is no easy way to check whether a given vertex $v$ is contained in the large clique. To prove the following theorem we can use that $\CK_k^{1/2}$ is isomorphism closed to translate the case for $\SK_k^{1/2}$ at least for the case of cliques of standard size.

\begin{thrm}\label{thrmck}
Let $c\in\N$ then 
\[K(\CK_k^{1/2},F,G)\bbl\text{$E$ has a clique of size $c$}\bbr>\0.\]
\end{thrm}
\begin{proof}
Consider the set \[\G_n^+=\{(G,U);\text{$G\in \G_n$ and $U\subseteq \set{n}$ is a $\floor{n/2}$ clique in $G$ }\}\]
and projections
\begin{align*}
\pi_1:\G_n^+&\to\G_n,(G,U)\mapsto G.\\
\pi_2:\G_n^+&\to\Po(\set{n}),(G,U)\mapsto U
\end{align*}

From the fact that $\CK_k^{1/2}$ is isomorphism closed, we have that $\abs{\pi_2^{-1}[U_1]}=\abs{\pi_2^{-1}[U_2]}$. We will set $\nu:=\abs{\pi_2^{-1}[U]}$ for some clique $U$ of size $\floor{n/2}$ so we have $\nu\abs{\G_n}=\G_n^+$. Let $\varphi_c(\overline x)$ be a $\{E\}$-sentence is satisfied iff $\overline x$ is a $c$ sized clique.  Denote an event in the sample space $\G_n$
\[A_{\overline a} :=\{G\in\G_n;G\models \varphi_c(\overline a)\}.\]
and in the sample space $\G_n^+$
\[A_{\overline a}^+ :=\{(G,U)\in\G_n^+;\overline a\text{ is a subclique of $U$}\}\]
notice that if $(G,U)\in A_{\overline a}^+$ then $G\in A_{\overline a}$ which implies (by another argument with projections) that
\[\nu\cdot\abs{A_{\overline a}}\geq \abs{A_{\overline a}^+}\]
so in conjunction with Example \ref{examskhalf} we have
\begin{align*}
\Pr_{G\in\G_n}[A_{\overline a}]&\geq \frac{\abs{A_{\overline a}^+}}{\abs{\G_n^+}}\\
&= \Pr_{G\in\SK_{k}^{1/2}}[\overline a \text{ is in the clique}]\\
&>(1/2)^c
\end{align*}
and since $c$ is constant we have $\Pr_{G\in\G_n}[A_{\overline a}]>0$.
\end{proof}

Of course it remains to show that using trees can actually increase the value all the way to $\1$, we did not get to prove that.

Finally one can return to cliques of nonstandard size and intuitively one expects this to be hard. For a clique of size greater than $n^{1/t}$ for any $t>\N$ one has to check more than $n^{1/t}$ edges to even know whether a given set is a clique and the counting argument used in the previous theorem implies that just guessing the clique of nonstandard size is not enough. We therefore present the following conjecture.

\begin{conj}
Let $m\leq c \ln n$ for any $c\in\N$, then
\[K(\CK_k^{1/2},F,G)\bbl\text{$E$ has a clique of size $\floor{n/m}$}\bbr=\0.\]
\end{conj}


\chapter*{Concluding remarks}
\addcontentsline{toc}{chapter}{Concluding remarks}

In this thesis we built a basic theory around wide limits of graphs, proved several general theorems and described the theories of the wide limits $\lim_F\EDGE_n$, $\lim_{F_{rud}}\ALL_n$, $\lim_{F_{nbtree}}\pPATH_n$, $\lim_{F_{nbtree}}\pPATHl_n$ and $\lim_{F_{dtree}}\pDPATH_n$ and proved that they are complete (Corollary \ref{crlltoosparse}, Theorem \ref{thrmall} and corollaries \ref{crllpPATH}, \ref{crllpPATHl} and \ref{crllpDPATH}).  We also proved that $\lim_{F_{rud}}^{G_{rud}}\SK_n^{1/2}$ contains a clique of size $\floor{n/(2\ln n)}$ (Theorem \ref{thrmsk}) and also that nonexistence of finite cliques in not valid in the more complex wide limit $\lim_{F_{rud}}^{G_{rud}}\CK_n^{1/2}$ (Theorem \ref{thrmck}).

During the development we planned to analyze wide limits the family $F_{poly}$ of polynomial functions. In \cite{krajicek2010forcing} it was proven that forcing with $F_{poly}$ results in quantifier elimination which implies that if an $\{E\}$-sentence holds in large enough $\G_k$, it has to hold in the limit. However the second order limit can still provide some information about the ability of polynomial time functions to search interesting subsets of $G\in\G_k$. In the end we did not get any new results about it. We want to mention that even though it would seem that the properties of $F_{poly}$ limits directly correspond to the $\P$ vs. $\NP$ question it seems that something a bit different happens. The way the limit objects are defined, it is not enough that some polynomial time algorithm does not exist to see that we cannot witness some property in the limit, but it is important that no polynomial algorithm succeeds on any standard fraction of inputs. This more closely resembles the generic case polynomial time (see \cite{gilman2007report}).

Another natural question would be to consider structures over general languages than just the language of graphs, and allow trees to query atomic sentences. Other combinatorial structures like hypergraphs and tournaments could be considered. Furthermore, wide limits of finite universal algebras could be considered which could require a whole new theory. This all leads to the fact that generalized spectra, elementary classes of $\Sigma_1^1$-logic (see \cite{Fagin74}), with restricted vertex sets make up a wide sequence, there could be a connection to the theory of spectra of sentences.


Considering structures over arbitrary vocabularies opens the door to connecting the concept of a wide limit with the complexity of the Constraint Satisfaction Problem. Solutions to $\text{CSP}(\A)$ form a wide sequence and depending on the tractability of the class we could expect different behaviour from $\lim_{F}\text{CSP}(\A)_n$.

Finally, important direction in studying wide sequences would be to characterize some $\lim_F\G_n$ without the direct construction and therefore to prove upper and lower bounds for search problems in $F$.


%%% Bibliography
%%% Bibliography (literature used as a source)
%%%
%%% We employ bibTeX to construct the bibliography. It processes
%%% citations in the text (e.g., the \cite{...} macro) and looks up
%%% relevant entries in the bibliography.bib file.
%%%
%%% The \bibliographystyle command selects, which style will be used
%%% for references from the text. The argument in curly brackets is
%%% the name of the corresponding style file (*.bst). Both styles
%%% mentioned in this template are included in LaTeX distributions.

\bibliographystyle{plainnat}    %% Author (year)
% \bibliographystyle{unsrt}     %% [number]

\renewcommand{\bibname}{Bibliography}

%%% Generate the bibliography. Beware that if you cited no works,
%%% the empty list will be omitted completely.

\bibliography{bibliography}

%%% If case you prefer to write the bibliography manually (without bibTeX),
%%% you can use the following. Please follow the ISO 690 standard and
%%% citation conventions of your field of research.

% \begin{thebibliography}{99}
%
% \bibitem{lamport94}
%   {\sc Lamport,} Leslie.
%   \emph{\LaTeX: A Document Preparation System}.
%   2nd edition.
%   Massachusetts: Addison Wesley, 1994.
%   ISBN 0-201-52983-1.
%
% \end{thebibliography}


%%% Figures used in the thesis (consider if this is needed)
\listoffigures

%%% Tables used in the thesis (consider if this is needed)
%%% In mathematical theses, it could be better to move the list of tables to the beginning of the thesis.
\listoftables

%%% Abbreviations used in the thesis, if any, including their explanation
%%% In mathematical theses, it could be better to move the list of abbreviations to the beginning of the thesis.
\chapwithtoc{List of Abbreviations}

%%% Attachments to the master thesis, if any. Each attachment must be
%%% referred to at least once from the text of the thesis. Attachments
%%% are numbered.
%%%
%%% The printed version should preferably contain attachments, which can be
%%% read (additional tables and charts, supplementary text, examples of
%%% program output, etc.). The electronic version is more suited for attachments
%%% which will likely be used in an electronic form rather than read (program
%%% source code, data files, interactive charts, etc.). Electronic attachments
%%% should be uploaded to SIS and optionally also included in the thesis on a~CD/DVD.
%%% Allowed file formats are specified in provision of the rector no. 72/2017.
\appendix
\chapter{Attachments}

\section{First Attachment}

\openright
\end{document}
